\documentclass[a4paper, 11pt]{article}
\usepackage[left=2.5cm, right=2.0cm, top=2.5cm, bottom=1.5cm, headheight=14pt]{geometry}
\usepackage[utf8]{inputenc}
\usepackage[english]{babel}
\usepackage{csquotes}
\usepackage{helvet}
\usepackage{titlesec}
\usepackage{fancyhdr}
\usepackage{xcolor}
\usepackage{lastpage}
\usepackage{enumitem}
\usepackage{setspace}
\usepackage{array}
\usepackage[colorlinks=true, linkcolor=blue, urlcolor=blue]{hyperref}
\usepackage[backend=biber, style=numeric-comp]{biblatex}

\renewcommand{\familydefault}{\sfdefault}
\setlength{\parindent}{0pt}
\linespread{1.15}
\setlength{\parskip}{2.75ex}
\setlist[itemize]{topsep=-1ex, leftmargin=1.5em}

\definecolor{CommentColor}{RGB}{140, 140, 140}

\titlespacing{\section}{0ex}{\parskip}{0ex}
\titlespacing{\subsection}{0ex}{\parskip}{0ex}
\titlespacing{\subsubsection}{0ex}{\parskip}{0ex}
\titlespacing{\paragraph}{0ex}{\parskip}{0ex}

\setcounter{secnumdepth}{4}
\titleformat{\section}{\normalfont}{\bf\thesection}{1em}{\bf}
\titleformat{\subsection}{\normalfont}{\bf\thesubsection}{1em}{\bf}
\titleformat{\subsubsection}{\normalfont}{\bf\thesubsubsection}{1em}{\bf}
\titleformat{\paragraph}{\normalfont}{\bf\theparagraph}{1em}{\bf}

\pagestyle{fancy}
\fancyhf{}
\fancyhead[L]{{\scriptsize Based on DFG form 3.06 -- 03/25}}
\fancyhead[R]{{\scriptsize page \thepage{} of \pageref*{LastPage}}}
\renewcommand{\headrulewidth}{0pt}

\newcommand{\heading}[1]{{\Large\textbf{#1}\par\bigskip}}
\newenvironment{commentblock}{\begin{spacing}{1.25}\color{CommentColor}}{\end{spacing}\bigskip}
\newcommand{\comment}[1]{{\color{CommentColor}#1}}
\newcommand{\largecomment}[1]{{\Large\textbf{\comment{#1}}}}
  
\newcommand{\tg}{\textgreater}
\newcommand{\tl}{\textless}
\newcommand{\horizontalrule}{\rule{\textwidth}{0.4pt}}

\addbibresource{abbrevs.bib}
\addbibresource{references.bib}


\begin{document}

\largecomment{\tl{}PUBLIC PART\tg}

\largecomment{\tl{}FILE\tg}

\begin{commentblock}
\tl{}This section is intended for reviewers and for the DFG's Head Office and statutory bodies,
but it may be published by the author of the report on a voluntary basis.\tg
\end{commentblock}


\heading{FINAL REPORT}


\section{General information}

DFG reference number: \comment{\tl{}Your DFG reference number\tg}

Project number: \comment{\tl{}You will find this in the letter of approval.\tg}

Project title:

Name(s) of the applicant(s):

Official address(es):

Name(s) of the co-applicants:

Names(s) of the cooperation partners:

Reporting period (entire funding period):

\newpage


\section{Summary}

\begin{commentblock}
\tl{}Provide a generally understandable summary in German and English (maximum 3,000 char-
acters each) in which you present the topic and describe the relevance of the project’s findings
to the interested public.\tg
\end{commentblock}


\section{Progress report}

\begin{commentblock}
\tl{}Address the following points in your report as applicable:

\begin{itemize}

\item Background and objectives of the project

\item Description of the project-specific results and findings. Please refer to the contributions of all applicants as well as other persons involved, cooperation partners, etc. Results that are already generally accessible in published form can be briefly summarised with refer- ence to the publication. Unpublished results are to be described in more detail.
  
\item Deviations from the original concept; findings that contradict the initial hypothesis
  
\item Activities and approaches to quality-enhancing measures through which the validity or verifiability of your research findings was ensured

\item Description of the handling of research data generated in the project and the data infrastructures used, if any (use the following checklist as a guide:
    \href{https://www.dfg.de/research\_data/checklist}{www.dfg.de/research\_data/checklist})

\item Description of any research data, methods, standards, software or infrastructures gener- ated in the project that are re-usable and openly accessible to others
    
\item Implementation of scientific events, science communication measures
  
\item Bibliography (list of works you referred to in describing the scientific results generated by the project and putting these in context. This might include your own work and that of other researchers.)\tg
  
\end{itemize}
\end{commentblock}


\section{Published Project Results}

\begin{commentblock}
\tl List here the main results that have emerged directly from the project and have been published; include the DOI (Digital Object Identifier), ISBN or other persistent identification number wherever possible. If this is not available, please provide the direct link. If the medium permits, publications must contain a reference to DFG funding (e.g. by means of a funding acknowledgement) and the project number. Structure the published project results as follows:\tg
\end{commentblock}


\subsection{Category A -- Articles in peer-reviewed journals, contributions to peer-reviewed
  conferences or to anthology volumes, and book publications}

\begin{commentblock}
\tl In this category please enter articles in peer-reviewed journals, contributions to peer-reviewed conferences or to anthology volumes, and book publications (see also \href{https://www.dfg.de/formulare/1_91}{DFG form 1.91}). Open access publications should be designated accordingly.\tg
\end{commentblock}


\subsection{Category B – Any other form of published results}

\begin{commentblock}
\tl Here you can cite any other form of published research results and findings. This may include contributions to non-peer-reviewed conferences, articles on preprint servers, data sets, protocols of clinical trials, software packages, blog contributions, infrastructures or transfer. You may also indicate other forms of academic output here, such as contributions to the (technical) infrastructure of an academic community (including in an international context) and contributions to science communication.\tg
\end{commentblock}


\subsection{Patents (applied for and granted)}

\newpage


\largecomment{\tl NON-PUBLIC PART\tg}

\section{Further information on the project, qualifications and outlook}

\begin{commentblock}
\tl This section is only intended for reviewers and for the DFG’s Head Office and statutory bodies; it is \underline{not intended for publication}.\tg
\end{commentblock}

\begin{commentblock}
\tl Address the following points in your report:

\begin{itemize}

\item Description of the progress of the project, including any problems encountered in its organisation or implementation.

\item Where applicable, potential follow-up studies or indication of potential applications, in particular with regard to knowledge transfer.

\item Qualification of researchers in early career phases in connection with the project (e.g. doctorates, post-doctoral lecturing qualifications, etc.). Use the table (5.1) to draw up a systematic list of the doctoral projects that were funded as part of the project.\tg

\end{itemize}
\end{commentblock}


\subsection{Doctoral researchers involved}

\newcolumntype{P}[1]{>{\raggedright\arraybackslash}p{#1}}
\begin{tabular}{|P{2.9cm}|P{2.0cm}|P{3.1cm}|P{3.15cm}|P{3.15cm}|}
\hline
Doctoral researchers
& Gender\newline (m/f/d)
& Doctoral status\newline (ongoing, finished, discontinued)
& Start (and where applicable) finish of doctoral studies\newline MM/YYYY -- MM/YYYY
& Funding within the framework of the project\newline MM/YYYY -- MM/YYYY
\\
\hline
Surname, first name
&
&
&
&
\\
\hline
Surname, first name
&
&
&
&
\\
\hline
\dots
&
&
&
&
\\
\hline
\end{tabular}

\end{document}


